\chapter{Формирование конкурентных вариантов схем сети, включая выбор номинального напряжения участков сети}
\label{cha:var_scheme}

\section{Разработка вариантов схем сети}

В рамках данного курсового проекта необходимо сформировать два варианта схем электрической сети включающих в себя как сети радиально-магистрального типа, так и кольцевую сеть.

Поскольку в составе потребителей всех пунктов имеется нагрузка I и II категории надежности, для обеспечения бесперебойности электроснабжения согласно пункту п. 1.2.19 и п. 1.2.20 ПУЭ \cite{пуэ7} требуется предусмотреть возможность получения электроэнергии от двух независимых взаимнорезервирующих источников питания. Это требование исключает возможность применения одноцепных линий в составе разомкнутых сетей радиально-магистрального типа. Таким образом, проектируемая районная электрическая сеть должна содержать участки двухцепных радиально-магистральных линий, либо сформированные в кольцевую сеть одноцепные линии.

Выбор среди конкурирующих вариантов осуществляется, опираясь на технические и технико-экономические показатели:
\begin{itemize}
	\item передача электроэнергии потребителям должна осуществляться по возможно кратчайшему пути, что обеспечивает снижение стоимости сооружения линий и экономию потерь мощности и электроэнергии;
	\item схема сети должна обеспечивать требуемый уровень надежности электроснабжения потребителей;
	\item схема сети должна быть по возможности (обснованно) простой;
	\item следует стремиться к минимизации количества трансформаций напряжения, что снижает необходимую установленную мощность трансформаторов и автотрансформаторов и, соответственно, капиталовложения на сооружение сети, а также - потери мощности и электроэнергии;
	\item не допускается сооружение линий по параллельным, рядом идущим трассам;
	\item следует избегать строительства малозагруженных линий, используемых только во время отключения элементов сети;
	\item не рекомендуется сооружать кольцевые сети, обеспечивающие электроснабжение 4-5 подстанций, из-за недопустимо больших потерь напряжения в послеаварийных режимах;
	\item комплекс номинального напряжения и схемы сети должны обеспечивать необходимое качество электроснабжения потребителей и выполнение технических ограничений по параметрам электрооборудования линий и подстанций;
	\item возможность сохранения принятых решений по развитию сети при небольших отклонениях нагрузок от прогнозируемых.
\end{itemize}

Изобразим возможные варианты схем районных сетей для наших исходных данных на рисунке \ref{fig:variant_scheme}.

Первый вариант [рисунок \ref{fig:variant_scheme} (а)] "--- кольцевая схема сети. Данный вариант соответствует требованиям по уровню надежности, обладает малым количеством числа трансформации уровней напряжения. Так же этот вариант обладает простыми схемами транзитных распределительных устройств (РУ). К недостатком такой схемы сети можно отнести повышенные потери мощности и электроэнергии и большими потерями напряжения в послеаварийных режимах, а также повышенная длина трасс линий.

Второй вариант [рисунок \ref{fig:variant_scheme} (б)] "--- радиально-магистральная схема сети с небольшой кольцевой сетью в конце. Такой вариант мы сразу можем отбросить, так как не допускается сооружение линий по рядом идущим трассам.

Третий вариант [рисунок \ref{fig:variant_scheme} (в)] "--- схема сети радиально-магистрального типа. Такие схемы обладают наименьшей длиной трасс линий, наименьшими потерями напряжения, мощности и электроэнергии и большими резервами по пропускной способности линий при перспективном росте нагрузок в заданных пунктах по сравнению с кольцевыми сетями. Однако этот вариант не рекомендуется к выбору, так как в нем соединяются все 5 ПС в одну линию, что из опыта проектирования ведет к тому, что на дальних пунктах тяжело поддерживается нужное напряжение.

Четвертый вариант [рисунок \ref{fig:variant_scheme} (г)] "--- является более оптимальной радиально-магистральной схемой сети, по сравнению с предыдущим вариантом, так как она лишена недостатка поддержания напряжения на последней ПС.

Таким образом, можем сделать вывод, что наиболее целесообразными конкурентными схемами, будут первая и четвертая (рисунки \ref{fig:variant_scheme} (а) и (г) соответственно), далее схемы №1 и №2.

\begin{figure}[h]
	\centering
	\includegraphics[width=0.7\textwidth]{inc/svg/variant_scheme}
	\caption{Варианты схем районных сетей}
	\label{fig:variant_scheme}
\end{figure}


\section{Выбор номинальных напряжений}

Экономически целесообразное номинальное напряжение участка сети зависит от некоторых параметров, среди которых основными являются передаваемая активная мощность по одной цепи линии и длина линии электропередачи.

\begin{table}[ht]
	\small
	\caption{Ориентировочные значения длин линий, дальности электропередачи и передаваемых мощностей для номинальных напряжений 35-220 кВ электрических сетей}
	\begin{tabularx}{\textwidth}{|X|Z|Z|Z|}
		\hline
		Номинальное напряжение, кВ & Экономически целесообразная передаваемая мощность на одну цепь линии, МВт & Средняя длина линий между соседними подстанциями, км & Средняя дальность электропередачи, км \\ \hline
		35 & 3 "--- 8 & 10 & 20 \\ \hline
		110 & 10 "--- 45 & 25 & 75 \\ \hline
		220 & 70 "--- 140 & 100 & 200 \\ \hline
	\end{tabularx}
	\label{tab:orient_l}
\end{table}

Выполним предварительный расчет потокораспределения активной мощности, передаваемой по одной цепи каждой линии, при котором допускается не учитывать потери мощности в сети. Расчет для схемы №1:
\[P_{4-5}^{\textup{1ц}} = \frac{P_5}{2} = \frac{35}{2} = 17,5\; \textup{МВт}\]
\[P_{3-5}^{\textup{1ц}} = \frac{P_4 + P_5}{2} = \frac{40 + 35}{2} = 37,5\; \textup{МВт}\]
	
Для предварительного расчета потокораспределения активной мощности в кольцевой сети К-1-3-2-К представим нагрузки пунктов 3-4-5 в виде нагрузки одного эквивалентного пункта 3':
\[P_3{'} = P_3 + P_4 + P_5 = 30 + 40 + 35 = 105\; \textup{МВт}\]

Вычислим активную мощность передаваемую по головным участкам кольцевой сети. Принимаем во внимание, что на этапе проектирования, когда сечение проводов ЛЭП еще не выбраны, для всех линий кольцевой сети, разрешается рассчитывать предварительное потокораспределение по длинам линии вместо сопротивлений.

Для участка K-1:
%\begin{eqndesc}[H]
\begin{equation*}
	\begin{split}
		P_\textup{К-1} &= \frac{P_1(L_{13} + L_{23} + L_\textup{К'-2}) + P_3{'}(L_{23} + L_\textup{К'-2}) + P_2\cdot L_\textup{К'-2}}{L_\textup{К-1} + L_{13} + L_{23} + L_\textup{К'-2}} = \\
			  &= \frac{70(49,2 + 41,8 + 49,2) + 105(41,8 + 49,2) + 70\cdot 49,2}{41,8 + 49,2 + 41,8 + 49,2} = 125,3\; \textup{МВт}
	\end{split}
\end{equation*}
%\end{eqndesc}

Для участка К'-2:
\begin{equation*}
	\begin{split}
		P_\textup{К'-2} &= \frac{P_2(L_{23} + L_{13} + L_\textup{К-1}) + P_3{'}(L_{13} + L_\textup{К-1}) + P_1\cdot L_\textup{K-1}}{L_\textup{К-1} + L_{13} + L_{23} + L_\textup{К'-2}} = \\
		&= \frac{70(41,8 + 49,2 + 41,8) + 105(49,2 + 41,8) + 70\cdot 41,8}{41,8 + 49,2 + 41,8 + 49,2} = 119,7\; \textup{МВт}
	\end{split}
\end{equation*}

Проверим правильность полученных результатов расчета активной мощности на головных участках кольцевой сети посредством оценки баланса мощности:
\[P_\textup{К-1} + P_\textup{К'-2} = 125,3 + 119,7 = 245\; \textup{МВт}\]
\[P_1 + P_3{'} + P_2 = 70 + 105 + 70 = 245\; \textup{МВт}\]
проверка сошлась

Теперь определим потокораспределение активной мощности на участках 1-3 и 2-3:
\[P_{13} = P_\textup{K-1} - P_1 = 125,3 - 70 = 55,3\; \textup{МВт}\]
\[P_{23} = P_\textup{K'-1} - P_2 = 119,7 - 70 = 49,7\; \textup{МВт}\]

Так как \(P_{13} > 0\) и \(P_{23} > 0\), то точка потокораздела находтся в третьем пункте потребления.

Полученные результаты предварительного потокораспределения активной мощности для удобства последующих рассуждений о выборе номинального напряжения сведем в табл. \ref{tab:vybor_u}

\begin{table}[H]
	\small
	\caption{Выбор номинального напряжения участков сети}
	\begin{tabularx}{\textwidth}{|l|Z|Z|Z|Z|Z|Z|}
		\hline
		Линия                  & K1    & K2    & 13   & 23   & 35      & 45   \\ \hline
		\(P_\textup{1ц}\) МВт  & 125,3 & 119,7 & 55,3 & 49,7 & 37,5    & 17,5 \\ \hline
		L, км                  & 41,8  & 49,2  & 49,2 & 41,8 & 46,4    & 25,9 \\ \hline
		\(U_\textup{ном}\), кВ & 220   & 220   & 220  & 220  & 110/220 & 110 \\ \hline
	\end{tabularx}
	\label{tab:vybor_u}
\end{table}

Все участки кольцевой сети следует проектировать на одно номинальное напряжение, выбираемое по наиболее загруженным головным участкам. Поскольку по головным участкам кольцевой сети К-1, К-2 длиной 41,8 и 49,2 км передается активная мощность более 70 МВт, в соответствии с таблицей \ref{tab:orient_l} для всех участков кольцевой сети следует выбрать номинальное напряжение 220 кВ.
%%% Local Variables:&
%%% mode: latex
%%% TeX-master: t
%%% End: